\documentclass[12pt, a4paper]{article} % Classe: article, report o book [8]

% --- Lingua e Codifica (Supporto Italiano) ---
\usepackage[utf8]{inputenc}  % Codifica input
\usepackage[T1]{fontenc}      % Codifica font
\usepackage[italian]{babel}   % Lingua italiana [2]

% --- Impaginazione e Margini ---
\usepackage[margin=2.5cm]{geometry} % Imposta margini (2.5cm su tutti i lati)

% --- Pacchetti Scientifici/Matematici ---
\usepackage{amsmath, amssymb, amsfonts, amsthm} % Matematica avanzata [2]
\usepackage{graphicx} % Per inserire immagini
\usepackage{hyperref} % Per collegamenti ipertestuali
\usepackage{booktabs} % Per tabelle professionali

% --- Formattazione Testo e Altro ---
\usepackage[document]{ragged2e}
\usepackage{xcolor}    % Per usare i colori 
\title{Alessandro Sassi}
\date{09-01-2026}


\begin{document}

\maketitle

\section{Ruolo}
Alessandro Sassi è un commerciale con esperienza di circa 6 mesi all'interno dell'azienda e si occupa di :
\begin{itemize}
      \item Creazione delle offerte all'interno di Access per conto di altri commerciali.
      \item Creazione di preventivi di richieste di assistenza.
\end{itemize}

\section{Operazioni comuni}
Le operazioni più comuni svolte da Alessandro all'interno dell'applicativo
includono:
\begin{itemize}
      \item Creazione di offerte per ricambi e assistenza (RA).
      \item Creazione delle offerte per nuovi impianti.
      \item Creazione degli ordini qualvolta volta un'offerta viene accettata.
      \item Nella creazione delle offerte ricerca e copia di vecchie offerte simili.
\end{itemize}

\section{Flusso di lavoro attaule}
Il flusso di lavoro attuale per la creazione delle offerte prevede i seguenti
passaggi:
\begin{enumerate}
      \item Arrivano dai commerciali più esperti o agenti degli schemi abbozzatti (anche
            collage di immagini) di quello che il cliente richiede.
      \item Viene aperto Access e si crea il cliente se manca.
      \item Si apre la schermata di creazione diversa a seconda del tipo di offerta
            (ricambi/assistenza o nuovi impianti).
      \item Si crea una nuova offerta copiando un ultima fatta per lo stesso cliente.
            \textcolor{red}{(Questo passaggio può portare errori in quanto la pagina non si
                  ricarca e si rischia di modificare l'offerta copiata di partenza)}.
      \item Si callenno a questo punto tutte le vecchie righe per partire a compilare da
            zero.
      \item Si inizia la compilazione dell'offerta:
            \begin{itemize}
                  \item Guardando lo schizzo del commerciale si inizia a cercare i componenti necessari
                        all'interno di access.
                  \item Si ricercano i codici solo in base alle descrizioni, ma l'unico modo per è
                        l'esperienza e gli appunti cartacei personali.
                  \item Certi componenti richiedendo obbligatoriamente la presenza di determinati
                        accessori (es. Globosilo $\rightarrow$ antiscoppio, cepatic, sonde, filtri
                        \dots) che vanno aggiunti manualmente. Questa infomazioni rimnagono
                        all'eperinza del singolo.
                  \item Alcuni componenti richiedono calcoli di costo esterni (es. pneumatici,
                        elettrici, dimensionali) che vanno fatti su file excel esterni e i risultati
                        inseriti manualmente.
                        I calcoli sono salvati e mantenuti in una dyrectory della commessa.
                  \item La proposta di un denterminata tipologia di silo è subbordinata all'utilizza ,
                        un silo per lo zuchero ha caratteristiche e prezzi diversi rispetto a un silo
                        per la farina. Qui dei template possono aiutare l'utente.
                  \item Note sulla funzionalalita possono aiutare l'utente meno asperto ad orientarsi.
                  Al momento esiste la schemarta "Prodotti".
                  \item Nell'iserimento bisogna stare attenti a non proporre compoennti non più di utilizzo ,
                  al momento non è presente un sistema che avverta di ciò a meno della presenza di "EX-" prima della descrizione.
                  \item Inserimento condizioni di montaggio, pagamento e consegna.
                  \item Inserimento dei costi e tempi di montaggio in base all'impianto ed il luogo di installazione.
            \end{itemize}
      \item Si compila riga per riga.
      \item Si ricerca i codice di access in base alla descrizione , se il prezzo non p
            rpesente ci si appoggia a file excel esterni.
      \item Si salva l'offerta e si manda al commerciale richiedente.
      \item Se l'offerta viene accettata si crea l'ordine in Access.
      \item Si manda l'ordine all'amministrazione per la creazione in Ad Hoc.
\end{enumerate}

\section{Situazione attuale}

\subsection{Criticità tecniche e interfaccia}
Attualmente, Alessandro riscontra diverse problematiche tecniche e di
interfaccia con l'applicativo Access:
\begin{itemize}
      \item \textbf{Obsolescenza della UI:} L'interfaccia attuale è percepita come ferma agli "anni '90", con un utilizzo inefficiente dello spazio a schermo e colori poco chiari .
      \item \textbf{Bug funzionali:}
            \begin{itemize}
                  \item Il sistema si blocca frequentemente , richiedendo attese o riavvii .
                  \item \textbf{Problema tastierino numerico:} Access disabilita il "Bloc Num" ogni volta che si aggiunge una riga o si stampa,
                        causando frequenti errori di digitazione nelle offerte .
            \end{itemize}
      \item \textbf{Gestione versionamento:} Quando si crea una nuova offerta (es. ampliamento o modifica), il numero offerta non si aggiorna visivamente in tempo reale,
            portando al rischio di modificare la versione precedente .
\end{itemize}

\subsection{Criticità operative}
Durante la creazione delle offerte, Alessandro affronta diverse sfide operative:
\begin{itemize}
      \item \textbf{Eccessiva manualità:} La compilazione delle offerte richiede la scrittura manuale riga per riga di descrizioni e codici (fino a 200+ righe),
            aumentando drasticamente i tempi e la possibilità di errore .
      \item \textbf{Curva di apprendimento:} Non essendoci procedure guidate, l'apprendimento è basato sulla memoria e su appunti personali/cartacei ("cheat sheets").
            Un junior impiega circa 4-6 mesi per diventare autonomo e ha ancora laucune sopperite dal cartaceo.
      \item \textbf{Frammentazione degli strumenti:} L'utente deve alternare continuamente tra il gestionale e numerosi file Excel esterni per calcoli (pneumatici, elettrici, dimensionali) .
      \item \textbf{Difficoltà di ricerca:} La ricerca dei componenti avviene solo per codice ed esperienza personale , nei ricambi se voglio un tipo preciso c'è margne di errore
            (Fatto l'esempio di una richiesta dove richiesto "Tutto ciò che è filtrante").
      \item \textbf{Compilazione manuale:} Molte riche di codice inserite manualmente senza allert o suggerimenti sono fonte di errrori.
      Se l'errore viene rilevato dopo la firma del cliente è difficilmente recuperabile a livello economico.
\end{itemize}

\section{Requisiti per il nuovo applicativo}

\subsection{Configurazione guidata}
Nel nuovo sistema, Alessandro richiede una serie di funzionalità per facilitare la creazione delle offerte:
\begin{itemize}
      \item \textbf{Selezione per accessori:} Si richiede un sistema che, selezionato un componente principale (es. "Globo"),
            proponga automaticamente gli accessori compatibili tramite checkbox o menu a tendina (es. sonde, filtri, valvole) .
      \item \textbf{Valori di Default Intelligenti:} Il sistema dovrebbe precompilare i campi standard (es. tubi diametro 100 per silo dello zucchero),
            riducendo l'input manuale .
      \item \textbf{Flessibilità obbligatoria:} È fondamentale che l'utente possa sempre sovrascrivere manualmente qualsiasi proposta del sistema (prezzi, descrizioni, componenti)
            per adattarsi a clienti o mercati specifici .
      \item \textbf{Contratti cartacei:} Ottimizzare i contratti cartacei a seconda del reparto di utilizzo , al momento sono presenti
            informazioni non utili e mancano altre utili (es. Dettagli di pagamento per reparti tecnici o dettagli tecnici per i reparti amministrativi).
\end{itemize}

\subsection{Integrazione e calcoli}
L'integrazione con altri sistemi e l'automazione dei calcoli sono aspetti chiave per migliorare l'efficienza:
\begin{itemize}
      \item \textbf{Visualizzazione 3D/dimensionale:} Integrazione desiderata con strumenti CAD (SolidWorks/Fusion)
            o logiche parametriche per fornire al cliente stime dimensionali preliminari in tempo reale durante la trattativa .
      \item \textbf{Integrazione calcoli ingegneristici:}
            \begin{itemize}
                  \item \textit{Calcolo Pneumatico:} Digitalizzazione dei grafici e delle logiche attualmente su Excel/Siemens per suggerire soffianti
                        e tubazioni in base alla portata e al batch .
                  \item \textit{Quadro Elettrico:} Deduzione automatica dei componenti del quadro in base agli elementi meccanici inseriti nell'offerta (es. numero di motori/stellari) .
                  \item \textit{Scale e Strutture:} Calcolo automatico del prezzo basato sui parametri del silo (altezza/diametro) .
            \end{itemize}
\end{itemize}

\subsection{Ricerca e organizzazione dati}
Per facilitare la selezione dei componenti, Alessandro propone le seguenti funzionalità:
\begin{itemize}
      \item \textbf{Ricerca per Tag/Argomento:} Implementazione di una ricerca semantica o per tag (es. "Safety") per raggruppare componenti correlati, oltre alla ricerca per codice .
      \item \textbf{Doppia Vista (vendita vs ricambi):} Possibilità di filtrare i cataloghi in base allo scopo:
            \begin{itemize}
                  \item \textbf{Vista vendita:} Componenti e accessori destinati alla vendita iniziale del macchinario .
                  \item \textbf{Vista ricambi:} Componenti di dettaglio per la manutenzione e la sostituzione post-vendita .
            \end{itemize}
      \item \textbf{Codifica unificata:} Necessità di allineare i codici tra gestionale, CAD (Fusion) e ERP (Ad Hoc) per facilitare il recupero delle informazioni .
\end{itemize} 

\section{Logiche di validazione e controllo}
Per minimizzare gli errori durante la creazione delle offerte, Alessandro suggerisce:
\begin{itemize}
      \item \textbf{Soft Alerts (Checklist):} Introduzione di un sistema di controllo finale che segnali eventuali dimenticanze (es. "Manca il vibrovaglio"), ma che \textbf{non sia bloccante}.
            L'utente deve poter confermare l'offerta anche in presenza di warning .
      \item \textbf{Prevenzione errori:} Evitare che modifiche successive al contratto o dimenticanze (es. tubazioni mancanti) generino costi non recuperabili una volta firmato l'ordine .
\end{itemize}

\section{Note Varie e Specifiche}
Alessandro ha evidenziato alcune esigenze specifiche da considerare:
\begin{itemize}
      \item \textbf{Incoterms:} Necessità di guidare la scelta tra i 19 Incoterms disponibili, fornendo descrizioni chiare per i commerciali meno esperti .
      \item \textbf{Localizzazione:} Gestione di prezzi, descrizioni e lingue variabili in base al paese di destinazione, mantenendo però la possibilità di deroga manuale .
      \item \textbf{Liquidi:} Necessità di digitalizzare le tabelle per la gestione dei componenti per liquidi (pompe, riscaldamento), attualmente gestite su vecchi supporti cartacei/Excel .
\end{itemize}

\section{Proposte future}
Alessandro ha proposto alcune idee per migliorare ulteriormente il processo
di vendita:
\begin{itemize}
      \item \textbf{Recoconto in tempo reale dei componenti suggeriti:} Un sistema che mostri in tempo reale il conteggio dei componenti
            suggeriti in base a quelli già selezionati rispetto ai modelli standard per tipologia impianto, aiutando l'utente a tenere traccia degli accessori necessari.
      \item \textbf{Divisione grafica per tipo di offerta:} Creare interfacce distinte per la creazione di offerte di ricambi/assistenza e nuovi impianti,
            ottimizzando l'esperienza utente in base al contesto specifico.
      \item \textbf{Visuliazzazione di modelli 3D:} Integrare un visualizzatore 3D all'interno dell'applicativo per permettere ai commerciali di
            mostrare ai clienti una rappresentazione visiva del macchinario proposto. Al momonto non possibile.
      \item \textbf{Proposte ricambistica:} Un sistema che suggerisca i ricambi comuni da proporre in base ai vecchi ordini.
\end{itemize}
\end{document}